\documentclass[11pt,a4paper]{article}

\usepackage[margin=1in]{geometry}
\usepackage{times}
\usepackage{amsmath,amssymb,amsfonts}
\usepackage{booktabs}
\usepackage{graphicx}
\usepackage{hyperref}
\usepackage{xcolor}
\usepackage{natbib}
\usepackage{url}
\usepackage{caption}

\hypersetup{
  colorlinks=true,
  linkcolor=blue!60!black,
  citecolor=blue!60!black,
  urlcolor=blue!60!black
}

\title{%
Quantifying Hamiltonian Inductive Bias under Symplectic Integration:\\
A Controlled Study on Conservative and Dissipative Pendulum Dynamics
}

\author{%
Jayacharan\\
\textit{Independent Researcher}\\
\texttt{https://github.com/Jayacharan09/pendulum-neural-dynamics}
}

\date{}

\begin{document}
\maketitle

\begin{abstract}
Neural networks trained solely on next-state prediction frequently achieve low single-step loss yet violate fundamental invariants over long time horizons. We present a tightly controlled empirical study of this failure mode on the simple pendulum. Two networks of matched capacity are trained on identical data: a standard feedforward map and a Hamiltonian Neural Network whose dynamics are obtained by automatic differentiation of a learned scalar Hamiltonian. Both models are integrated with the velocity Verlet method (a second-order symplectic integrator).

On the energy-conserving pendulum the Hamiltonian network reduces mean relative energy error from 26.5\% to 0.22\% over an 8-second rollout---an improvement of approximately two orders of magnitude. When the same architecture is applied to a damped pendulum (energy no longer conserved), the Hamiltonian inductive bias loses its advantage and slightly underperforms the unconstrained baseline. These results demonstrate that the benefit of Hamiltonian structure is conditional on the underlying physics matching the assumed geometric constraints, and that the combination of a correct inductive bias with a symplectic integrator yields substantially better long-horizon fidelity than either component alone.

\vspace{0.5em}
\noindent\textbf{Keywords:} Hamiltonian Neural Networks, symplectic integrators, energy conservation, inductive bias, physics-informed machine learning, long-horizon prediction
\end{abstract}

\section{Introduction}

Data-driven models of dynamical systems are attractive because they do not require explicit derivation of governing equations. In practice, however, networks that minimise single-step prediction error often produce trajectories that gradually violate conservation laws. Energy drift is especially destructive: small local errors accumulate and eventually destroy the qualitative correctness of the predicted motion.

The simple pendulum remains an ideal testbed. Its equations are elementary, its phase space is two-dimensional, and mechanical energy is exactly conserved in the frictionless case. Any learned model that fails to respect this invariant produces an unambiguous, visually clear failure mode (inward or outward spiralling in phase space).

Prior work introduced Hamiltonian Neural Networks (HNNs) that learn a scalar energy function and recover dynamics via Hamilton's equations. Most published evaluations, however, still rely on Euler integration, which itself introduces artificial energy drift. In this paper we isolate three factors under matched conditions:

\begin{enumerate}
  \item the presence or absence of Hamiltonian inductive bias,
  \item the use of a symplectic (velocity Verlet) integrator instead of Euler,
  \item the match or mismatch between the assumed physics and the true system (conservative versus damped pendulum).
\end{enumerate}

Our contribution is a clean quantitative measurement of residual energy error under these controlled variations.

\section{Related Work}

Hamiltonian Neural Networks were introduced by \citet{greydanus2019hamiltonian}. Subsequent work explored Lagrangian Neural Networks~\citep{lutter2019deep}, Hamiltonian Graph Networks~\citep{toth2019hamiltonian}, and symbolic regression with inductive biases~\citep{cranmer2020discovering}.

The present study is deliberately narrower. We do not propose a new architecture. We measure, under identical capacity, data, optimiser and integrator, how large the practical benefit of the Hamiltonian bias is, and under what conditions that benefit disappears.

\section{Method}

\subsection{Ground-truth dynamics}

The pendulum is governed by
\begin{equation}
  \ddot{\theta} = -\frac{g}{L}\sin\theta - c\,\dot{\theta},
\end{equation}
where \(c=0\) recovers the conservative case and \(c=0.15\) introduces linear damping. Ground-truth trajectories are generated with classical fourth-order Runge--Kutta.

\subsection{Networks}

Both networks use two hidden layers of 48 units with \(\tanh\) activations.

\begin{itemize}
  \item \textbf{Standard Network} learns the direct map \((\theta,\omega)\mapsto(\theta',\omega')\).
  \item \textbf{Hamiltonian Network} learns a scalar \(H(\theta,p)\) and obtains the vector field by automatic differentiation of Hamilton's equations:
  \begin{equation}
    \dot{\theta}=\frac{\partial H}{\partial p},\qquad
    \dot{p}=-\frac{\partial H}{\partial\theta}.
  \end{equation}
\end{itemize}

\subsection{Integrator}

All learned dynamics are advanced with the velocity Verlet method, a second-order symplectic integrator:
\begin{align}
  p_{n+1/2} &= p_n - \tfrac12\Delta t\,\partial_q H, \\
  q_{n+1}   &= q_n + \Delta t\,\partial_p H(q_n,p_{n+1/2}), \\
  p_{n+1}   &= p_{n+1/2} - \tfrac12\Delta t\,\partial_q H(q_{n+1},p_{n+1/2}).
\end{align}
This choice eliminates the dominant source of numerical energy drift that Euler integration would introduce.

\subsection{Training protocol}

\begin{itemize}
  \item Data: 60 random trajectories of length 8\,s (\(\Delta t=0.01\)).
  \item Optimiser: Adam, learning rate \(10^{-3}\), 25 epochs, batch size 256.
  \item Loss: mean-squared error on consecutive states.
  \item Evaluation: single unseen initial condition \(\theta_0=55^\circ\), \(\omega_0=0\).
\end{itemize}

\section{Results}

\subsection{Conservative pendulum (positive control)}

\begin{table}[h]
\centering
\begin{tabular}{lcc}
\toprule
Model & Mean relative energy error & Max relative energy error \\
\midrule
Standard Network           & 26.54\% & 37.50\% \\
Hamiltonian + Verlet       & \textbf{0.22\%} & \textbf{0.64\%} \\
\bottomrule
\end{tabular}
\caption{Energy conservation performance on the conservative pendulum (8-second rollout).}
\label{tab:cons}
\end{table}

The Hamiltonian network reduces mean energy error by a factor of approximately 120.

\subsection{Damped pendulum (negative control)}

\begin{table}[h]
\centering
\begin{tabular}{lcc}
\toprule
Model & Mean relative energy error & Max relative energy error \\
\midrule
Standard Network           & 37.09\% & 80.10\% \\
Hamiltonian + Verlet       & 40.92\% & 71.65\% \\
\bottomrule
\end{tabular}
\caption{Energy conservation performance on the damped pendulum (\(c=0.15\)). The Hamiltonian inductive bias loses its advantage.}
\label{tab:damped}
\end{table}

When energy is no longer a conserved quantity, the Hamiltonian inductive bias loses its advantage and slightly underperforms the unconstrained network.

\section{Discussion}

Under matched capacity, data and optimisation, replacing a direct state-transition map by a learned Hamiltonian and integrating it with a symplectic method reduces residual energy error by roughly two orders of magnitude on a conservative system. The same architectural choice becomes neutral or slightly harmful when the system is dissipative.

Two practical implications follow:

\begin{enumerate}
  \item Inductive bias is conditional. Encoding Hamiltonian structure is powerful only when the target system is (approximately) energy-conserving.
  \item Integrator choice matters. A non-negligible fraction of previously observed drift in earlier HNN evaluations was numerical rather than model-based.
\end{enumerate}

\section{Conclusion}

We have quantified the practical value of Hamiltonian inductive bias when paired with a symplectic integrator. On the energy-conserving pendulum the combination reduces long-horizon relative energy error from 26.5\% to 0.22\%. On a damped pendulum the same bias ceases to be beneficial. These controlled measurements confirm that geometric structure, when correctly matched to the underlying physics and integrated appropriately, remains one of the most effective tools for long-horizon prediction of dynamical systems.

Complete experimental code and results are available in the accompanying repository.

\bibliographystyle{plainnat}
\begin{thebibliography}{99}

\bibitem[Cranmer et al.(2020)]{cranmer2020discovering}
Miles Cranmer, Alvaro Sanchez-Gonzalez, Peter Battaglia, Rui Xu, Kyle Cranmer, David Spergel, and Shirley Ho.
\newblock Discovering symbolic models from deep learning with inductive biases.
\newblock In \emph{Advances in Neural Information Processing Systems}, 2020.

\bibitem[Greydanus et al.(2019)]{greydanus2019hamiltonian}
Samuel Greydanus, Misko Dzamba, and Jason Yosinski.
\newblock Hamiltonian neural networks.
\newblock In \emph{Advances in Neural Information Processing Systems}, 2019.

\bibitem[Hairer et al.(2006)]{hairer2006geometric}
Ernst Hairer, Christian Lubich, and Gerhard Wanner.
\newblock \emph{Geometric Numerical Integration}.
\newblock Springer, 2nd edition, 2006.

\bibitem[Lutter et al.(2019)]{lutter2019deep}
Michael Lutter, Christian Ritter, and Jan Peters.
\newblock Deep Lagrangian networks.
\newblock In \emph{International Conference on Learning Representations}, 2019.

\bibitem[Toth et al.(2019)]{toth2019hamiltonian}
Peter Toth et al.
\newblock Hamiltonian generative networks.
\newblock \emph{arXiv preprint arXiv:1909.13789}, 2019.

\end{thebibliography}

\end{document}
